\chapter{Introduction}
Integrating a new board into QEMU is a complex but essential process for emulating embedded hardware in a virtual environment. This project documents the addition of the NXP S32K3X8EVB board, based on the S32K358 System-on-Chip (SoC), into the QEMU ecosystem. The NXP S32K3X8EVB board features an ARM Cortex-M7 CPU and includes integrated peripherals such as LPUART (Low Power Universal Asynchronous Receiver-Transmitter) and PIT (Periodic Interrupt Timer).

This work follows several key phases:
\begin{enumerate}
\item Board Definition: Initial board configuration, including the implementation of the board initialization file, definition of system clocks, and creation of the fixed-frequency clock (SYSCLK) set to 48 MHz.
\item SoC Implementation: Creation of the SoC structure, including memory regions (such as flash memory), peripherals (including LPUART and PIT), and IRQ (Interrupt Request) lines for handling interrupts from the peripherals.
\item QEMU Integration: Modifying QEMU configuration files to include the new board and its corresponding SoC, including Kconfig entries for enabling the board in ARM architectures.
\item Peripheral Configuration: Detailed setup of LPUART modules for serial communication and PIT for time-based operations such as periodic interrupts and task scheduling.
\end{enumerate}

Through this process, it has been possible to accurately simulate the NXP S32K3X8EVB board within QEMU.
 

